% =====================================================================
% บทที่ 4 การใช้งานสำหรับนักศึกษา (คู่มือผู้ใช้งานทั่วไป: หนึ่งบทต่อหนึ่งบทบาท)
% =====================================================================
\label{ch:student}

\begin{guide}[แนวทางการเขียนคู่มือผู้ใช้งานทั่วไป]
เขียนหนึ่งบทต่อหนึ่งบทบาท (เช่น นักศึกษา อาจารย์ ผู้ดูแลรายวิชา) ถ้าระบบมีผู้ใช้บทบาทเดียวให้รวมเป็นบทเดียว
\begin{itemize}
  \item หนึ่งหัวข้อคือหนึ่งงาน ชื่อหัวข้อขึ้นต้นด้วย ``การ...'' เรียงตามลำดับที่ผู้ใช้ทำจริง
        ผู้อ่านจะค้นงานที่ต้องการจากสารบัญได้ทันที
  \item ทุกงานเขียนเป็นขั้นตอนด้วย \cmd{steps} และบอกผลลัพธ์ที่ควรเห็นด้วย \cmd{\result}
        ใส่ภาพหน้าจอเมื่อเป็นหน้าจอใหม่หรือหน้าจอที่มีหลายส่วน
  \item สถานะ ข้อความแจ้งเตือน และเงื่อนไข (เช่น แก้ไขได้เมื่อใด) เขียนเป็นตาราง
  \item กล่อง tip/note/warning/important ใช้เท่าที่จำเป็น ไม่เกิน 1--2 กล่องต่อหัวข้อ
  \item ไม่อธิบายการทำงานภายในของระบบ (เก็บไว้ในคู่มือสำหรับนักพัฒนา)
\end{itemize}
\end{guide}

\section{การสร้างข้อเสนอโครงงาน}
\label{sec:proposal}

นักศึกษาหัวหน้ากลุ่มเป็นผู้สร้างข้อเสนอโครงงาน แบบฟอร์มแบ่งเป็น 3 ส่วนดังภาพที่~\ref{fig:proposal}

\begin{screenshot}[width=0.5\linewidth]{images/proposal-form.png}{แบบฟอร์มรายละเอียดข้อเสนอโครงงาน}{fig:proposal}
  \callout{0.27}{0.07}{1}
  \callout{0.92}{0.237}{2}
  \callout{0.92}{0.60}{3}
  \callout{0.92}{0.76}{4}
\end{screenshot}

\begin{calloutlist}
  \item แถบสถานะของข้อเสนอ (สร้าง ส่งแล้ว รอพิจารณา)
  \item ส่วนที่ 1 ข้อมูลโครงงาน: ชื่อโครงงานภาษาไทย (จำเป็น) ชื่อภาษาอังกฤษ และรายละเอียด
  \item ส่วนที่ 2 สมาชิกโครงงาน: ผู้สร้างเป็นหัวหน้าโครงงานโดยอัตโนมัติ
  \item ส่วนที่ 3 อาจารย์ที่ปรึกษาหลัก (จำเป็น) และอาจารย์ที่ปรึกษาร่วม (ถ้ามี)
\end{calloutlist}

\begin{steps}
  \item ที่หน้ารายการโครงงาน คลิก \btn{+ สร้างข้อเสนอ}
  \item กรอก \field{ชื่อโครงงาน (ภาษาไทย)} \field{ชื่อโครงงาน (ภาษาอังกฤษ)}
        และ \field{รายละเอียด/ที่มาและความสำคัญ}
  \item เพิ่มสมาชิกโดยค้นหาจากชื่อ รหัสนักศึกษา หรืออีเมล (รวมหัวหน้าได้สูงสุด 3 คน)
  \item เลือก \field{อาจารย์ที่ปรึกษาหลัก} และ \field{อาจารย์ที่ปรึกษาร่วม} ถ้ามี
  \item คลิก \btn{บันทึกฉบับร่าง} เพื่อแก้ไขภายหลัง หรือ \btn{ส่งข้อเสนอ} เพื่อส่งให้อาจารย์พิจารณา
        \result{ข้อเสนอมีสถานะ \cmd{submitted} และอาจารย์ที่ปรึกษาได้รับการแจ้งเตือน}
\end{steps}

\begin{warning}
เลือกอาจารย์คนเดียวกันเป็นทั้งอาจารย์ที่ปรึกษาหลักและอาจารย์ที่ปรึกษาร่วมไม่ได้
และเมื่อส่งข้อเสนอแล้วจะแก้ไขไม่ได้จนกว่าอาจารย์จะส่งกลับให้แก้ไข
\end{warning}

\section{การติดตามสถานะข้อเสนอ}

ข้อเสนอโครงงานมี 4 สถานะ ตามตารางที่~\ref{tab:proposal-status}
เมื่อข้อเสนอได้รับอนุมัติ ระบบจะสร้างโครงงานพร้อมไฟล์แม่แบบ \LaTeX{} ให้โดยอัตโนมัติ

\begin{table}[H]
\caption{สถานะของข้อเสนอโครงงาน}
\label{tab:proposal-status}
\centering
\begin{tabular}{|L{2.6cm}|L{3cm}|L{7.6cm}|}
\hline
\textbf{สถานะในระบบ} & \textbf{ความหมาย} & \textbf{สิ่งที่นักศึกษาทำได้} \\ \hline
\cmd{draft}     & ฉบับร่าง       & แก้ไขและส่งข้อเสนอ \\ \hline
\cmd{submitted} & รอพิจารณา     & รอผลการพิจารณาจากอาจารย์ \\ \hline
\cmd{approved}  & อนุมัติแล้ว    & เริ่มเขียนเอกสารในโครงงานที่ระบบสร้างให้ \\ \hline
\cmd{rejected}  & ส่งกลับแก้ไข   & อ่านความเห็นของอาจารย์ แก้ไข แล้วส่งใหม่ \\ \hline
\end{tabular}
\end{table}

\section{การแก้ไขเอกสารด้วย LaTeX Editor}
\label{sec:editor}

คลิกการ์ดโครงงานที่อนุมัติแล้วเพื่อเปิด LaTeX Editor ซึ่งมีส่วนประกอบดังภาพที่~\ref{fig:editor}

\begin{screenshot}{images/editor.png}{ส่วนประกอบของหน้า LaTeX Editor}{fig:editor}
  \callout{0.53}{0.083}{1}
  \callout{0.062}{0.26}{2}
  \callout{0.34}{0.36}{3}
  \callout{0.78}{0.33}{4}
  \calloutline{0.817}{0.165}{0.817}{0.05}{5}
  \calloutline{0.863}{0.22}{0.863}{0.05}{6}
  \calloutline{0.908}{0.165}{0.908}{0.05}{7}
  \calloutline{0.963}{0.22}{0.963}{0.05}{8}
  \callout{0.22}{0.845}{9}
\end{screenshot}

\begin{calloutlist}
  \item แถบเครื่องมือ \LaTeX{} สำหรับแทรกบท หัวข้อ ตัวหนา ตาราง ภาพ และแผนภาพ Mermaid
  \item แฟ้มข้อมูลของโครงงาน สร้าง อัปโหลด และเปิดไฟล์
  \item พื้นที่แก้ไขโค้ด \LaTeX
  \item ตัวอย่าง PDF หลัง Compile
  \item ปุ่ม \btn{ซ่อน PDF} เพื่อขยายพื้นที่แก้ไข
  \item ปุ่มเปิดแผงความเห็น (หัวข้อ~\ref{sec:comment})
  \item ปุ่ม \btn{บันทึก}
  \item ปุ่ม \btn{Compile} สร้างไฟล์ PDF
  \item ผลการ Compile (COMPILE LOG) แสดงข้อผิดพลาดถ้ามี
\end{calloutlist}

แป้นพิมพ์ลัดที่ใช้บ่อยใน LaTeX Editor แสดงดังตารางที่~\ref{tab:shortcuts}

\begin{table}[H]
\caption{แป้นพิมพ์ลัดใน LaTeX Editor}
\label{tab:shortcuts}
\centering
\begin{tabular}{|L{4.6cm}|L{8.6cm}|}
\hline
\textbf{แป้นพิมพ์ลัด} & \textbf{การทำงาน} \\ \hline
\keys{Ctrl+S}       & บันทึกไฟล์ \\ \hline
\keys{Ctrl+Enter}   & Compile เอกสาร \\ \hline
\keys{Ctrl+Z}       & ย้อนกลับ (Undo) \\ \hline
\keys{Ctrl+Shift+Z} & ทำซ้ำ (Redo) \\ \hline
\keys{Ctrl+F}       & ค้นหา \\ \hline
\keys{Ctrl+H}       & ค้นหาและแทนที่ \\ \hline
\end{tabular}
\end{table}

\begin{tip}
เลือกข้อความก่อนกดปุ่ม \btn{ตัวหนา} หรือ \btn{ตัวเอียง} ระบบจะครอบข้อความด้วยคำสั่ง \LaTeX{} ให้
ถ้าไม่ได้เลือกข้อความ ระบบจะแทรกคำสั่งตัวอย่างที่ตำแหน่งเคอร์เซอร์
\end{tip}

\section{การแก้ไขเอกสารร่วมกับสมาชิก}

สมาชิกทุกคนในโครงงานเปิดไฟล์เดียวกันและแก้ไขพร้อมกันได้แบบเรียลไทม์
สถานะการเชื่อมต่อดูได้จากจุดสีที่แถบเครื่องมือ ตามตารางที่~\ref{tab:connection}

\begin{table}[H]
\caption{ความหมายของสถานะการเชื่อมต่อ}
\label{tab:connection}
\centering
% จุดสีวาดด้วย \tikz (ไม่ใช้ $\bullet$ ซึ่งระบบตรวจรูปเล่มนับเป็นสัญลักษณ์รายการ)
\begin{tabular}{|C{2.4cm}|L{10.8cm}|}
\hline
\textbf{สถานะ} & \textbf{ความหมายและสิ่งที่ควรทำ} \\ \hline
\tikz[baseline=-0.3em]\fill[cskkuTip] (0,0) circle (0.28em);~เขียว & เชื่อมต่อเรียลไทม์สำเร็จ เห็นการแก้ไขของสมาชิกคนอื่นทันที \\ \hline
\tikz[baseline=-0.3em]\fill[cskkuImp] (0,0) circle (0.28em);~แดง & ไม่ได้เชื่อมต่อ ทำงานแบบออฟไลน์และต้องบันทึกเอง \\ \hline
\tikz[baseline=-0.3em]\fill[cskkuWarn] (0,0) circle (0.28em);~เหลือง & เชื่อมต่อผิดพลาด ให้รีเฟรชหน้าเว็บ \\ \hline
\end{tabular}
\end{table}

\section{การ Compile และตรวจสอบไฟล์ PDF}

\begin{steps}
  \item คลิก \btn{Compile} หรือกด \keys{Ctrl+Enter}
  \item ตรวจผลลัพธ์ในพื้นที่ตัวอย่าง PDF (หมายเลข~\calloutmark{4} ในภาพที่~\ref{fig:editor})
  \item ถ้ามีข้อผิดพลาด อ่านรายละเอียดใน COMPILE LOG (หมายเลข~\calloutmark{9}) แล้วแก้ไขตามบรรทัดที่ระบุ
\end{steps}

\begin{note}
ระบบอนุญาตให้ Compile ได้ครั้งละ 1 คนต่อโครงงาน ถ้าเห็นข้อความ ``กำลังแปลผลอยู่''
แสดงว่าสมาชิกคนอื่นกำลัง Compile ให้รอสักครู่แล้วลองใหม่
\end{note}

\section{การส่งงาน}
\label{sec:submit}

\begin{steps}
  \item คลิกเมนู \btn{ส่งงาน} ที่แถบด้านบน
  \item คลิก \btn{สร้างการส่งงาน}
  \item เลือกโครงงานและกรอกหมายเหตุ
  \item ยืนยันการส่ง
        \result{ระบบบันทึกสำเนาของไฟล์ทั้งหมดในขณะนั้น (snapshot) พร้อมวันเวลาที่ส่ง}
\end{steps}

\begin{important}
ตรวจสอบกำหนดส่งที่หน้าแรกก่อนส่งทุกครั้ง การแก้ไขไฟล์หลังส่งจะไม่เปลี่ยนสำเนาที่ส่งไปแล้ว
\end{important}
